% from ega4-2-fr.tex
\oldpage[IV-2]{53}

\begin{enumerate}
\item[(i)] Si $f$ est quasi-compact et dominant, on a
$\dim(Y)\leqslant\dim(X)$.
\item[(ii)] Si $f$ est quasi-fini, on a
$\dim(X)\leqslant\dim(Y)$.
\item[(iii)] Supposons que $X$ soit de type fini sur $k$. Pour que
$\dim(X)\geqslant n$ (resp. $\dim(X)\leqslant n$, $\dim(X)=n$), il faut et il
suffit qu'il existe un ouvert $U$ dense dans $X$, et un $k$-morphisme
\[
g:U\longrightarrow\mathbf V(k^n)
=\operatorname{Spec}(k[T_1,\ldots,T_n]),
\]
que nous noterons aussi $\mathbf V_k^n$, qui soit surjectif (resp. fini,
resp. fini et surjectif). Si $X$ est un schéma affine, on peut prendre $U=X$.
\end{enumerate}
\end{theorem}

\begin{proof}
(i) Soit $y$ un point maximal de $Y$~; on sait que $f^{-1}(y)$ contient un
point maximal $x$ de $X$ (1.1.5), donc $k(x)$ est une extension de $k(y)$,
d'où l'inégalité
\[
\operatorname{deg.tr.}_k k(y)
\leqslant\operatorname{deg.tr.}_k k(x)
\leqslant\dim(X),
\]
ce qui démontre (i).

(ii) Si $x$ est point maximal de $X$, $k(x)$ est une extension finie de
$k(f(x))$ (II, 6.2.2), donc a même degré de transcendance sur $k$. Si l'on
considère le sous-préschéma fermé réduit de $Y$ ayant $\overline{\{f(x)\}}$
comme espace sous-jacent, on voit qu'on est ramené à démontrer le corollaire
suivant.
\end{proof}

\begin{corollary}[4.1.2.1]
\label{IV.2.4.1.2.1-fr}
Soient $Y$ un $k$-préschéma localement de type fini~; pour tout sous-préschéma
$Z$ de $Y$, on a $\dim(Z)\leqslant\dim(Y)$. Supposons en outre que les
composantes irréductibles de $Y$ aient toutes même dimension~; alors, pour que
$Z$ soit rare dans $Y$, il faut et il suffit que
$\dim(Z)<\dim(Y)$.
\end{corollary}

\begin{proof}
En effet, soit $z$ un point maximal de $Z$~; en considérant un des
sous-préschémas fermés réduits de $Y$ ayant pour espace sous-jacent une des
composantes irréductibles de $Y$ contenant $z$, on se ramène au cas où $Y$ est
intègre de point générique $y$, puis, en considérant dans $Y$ un ouvert affine
contenant $z$, au cas où $Y$ est affine et de type fini sur $k$, soit
$Y=\operatorname{Spec}(A)$, et $Z$ fermé dans $Y$, soit
$Z=\operatorname{Spec}(A/\mathfrak a)$ où $\mathfrak a$ est un idéal de
l'anneau intègre $A$, distinct de $A$. En vertu du lemme de normalisation
(Bourbaki, \emph{Alg. comm.}, chap.~V, \S~3,
n\textsuperscript{o}~1, th.~1), il existe une suite finie
$(x_i)_{1\leqslant i\leqslant n}$ d'éléments de $A$, algébriquement
indépendants sur $k$, tels que $A$ soit entier sur l'anneau
$B=k[x_1,\ldots,x_n]$ et que $\mathfrak a\cap B$ soit engendré par une
sous-famille $(x_i)_{1\leqslant i\leqslant p}$ (éventuellement vide) de
$(x_i)_{1\leqslant i\leqslant n}$. Soit
$g:Y\to\operatorname{Spec}(B)=\mathbf V_k^n$ le morphisme fini dominant qui
correspond à l'injection canonique $B\to A$~; comme
$C=B/(\mathfrak a\cap B)$ est isomorphe à $k[x_{p+1},\ldots,x_n]$, $g$ induit
sur $Z$ un morphisme fini dominant $Z\to\mathbf V_k^{n-p}$~; on a donc
$\operatorname{deg.tr.}_k k(y)=n$ et
$\operatorname{deg.tr.}_k k(z)=n-p$, puisque $k(y)$ (resp. $k(z)$) est une
extension finie de $k(g(y))$ (resp. $k(g(z))$). Cela démontre la première
assertion de (4.1.2.1). En outre, si $p=0$, on a nécessairement $z=y$, car $y$
est le seul point de $Y$ dont l'image par $g$ soit le point générique de
$\mathbf V_k^n$~; dans ce cas, $Z$ contient donc un ouvert non vide de $Y$. Si
au contraire $p>0$, on a nécessairement $z\neq y$ et
$\overline{\{z\}}$ est donc rare dans $\overline{\{y\}}$, ce qui achève de
prouver (4.1.2.1).
\end{proof}

\begin{proof}
(iii) Le fait que les conditions énoncées sont suffisantes résulte aussitôt de
(i) et (ii) et de (1.5.4, (v)). Pour prouver qu'elles sont nécessaires, on peut
considérer un ouvert dense $U$ de $X$, réunion d'ouverts affines $U_i$ deux à
deux disjoints et irréductibles $(1\leqslant i\leqslant m)$, dont chacun
contient un des points maximaux de $X$ $(0_{\mathrm I},2.1.6)$~; on peut en
outre supposer $X$ réduit et si $X$ est affine, on peut bien entendu prendre
$U=X$. Le même

